\chapter*{ВВЕДЕНИЕ}
\addcontentsline{toc}{chapter}{ВВЕДЕНИЕ}

\section*{Актуальность исследования}

Системная динамика предоставляет аппарат для исследования запасов, потоков и обратных связей в социальных, экономических, экологических и технических системах. Её практическое применение требует не только знания предметной области, но и умения перейти от качественной гипотезы к диаграммам причинно-следственных связей и потоков, а затем к размерно согласованной системе дифференциальных уравнений. Ошибка в направлении потока, единицах коэффициента или начальном состоянии может сделать формально исполняемую модель содержательно неверной.

Большие языковые модели способны извлекать переменные из естественно-языкового описания и предлагать формальные выражения. Однако вероятностная природа такой генерации приводит к пропускам переменных, нарушениям схемы, несовместимым единицам и меняющейся топологии. Синтаксически правильный JSON или успешно завершённая симуляция сами по себе не доказывают корректность модели. Возникает задача ограничить пространство допустимых результатов LLM так, чтобы сохранить гибкость текстового ввода и одновременно сделать проверку модели воспроизводимой.

В работе рассматривается подход, при котором LLM используется для интерпретации описания и планирования, а критические свойства результата обеспечиваются типизированным промежуточным представлением, детерминированными построителями и проверками. Это позволяет отделить ошибки понимания предметной области от ошибок представления и вычислений, а также сформулировать измеримые критерии качества.

\section*{Цель и задачи исследования}

\textbf{Целью работы} является разработка методов и алгоритмов автоматизированной генерации системно-динамических моделей на основе текстовых описаний с использованием больших языковых моделей и детерминированных ограничений, а также создание программной платформы, реализующей предложенный подход.

Для достижения цели поставлены следующие задачи:

\begin{enumerate}
    \item Проанализировать процесс перехода от текстовой гипотезы к исполняемой системно-динамической модели и существующие работы по применению LLM в моделировании.
    \item Разработать каноническое типизированное представление переменных, выражений, единиц измерения и топологии модели.
    \item Предложить конвейер, включающий краткое описание задачи, типизированный план компонентов, детерминированные построители, ограниченное исправление и контрольную симуляцию.
    \item Реализовать проверку единиц измерения, ссылок, направления потоков, согласованности CLD/SFD и поведенческих свойств результата.
    \item Реализовать CLI, REST API, backend и веб-интерфейс для контекстной работы с моделями, сценариями, анализом чувствительности и калибровкой.
    \item Обеспечить ограниченный импорт существующих Vensim-моделей и экспорт в Python, Vensim, XMILE и CSV.
    \item Провести инженерную и исследовательскую оценку реализованного подхода, зафиксировав как успешные, так и отрицательные результаты, а также границы их интерпретации.
\end{enumerate}

\section*{Объект и предмет исследования}

\textbf{Объектом исследования} являются процессы создания и проверки системно-динамических моделей с использованием методов искусственного интеллекта.

\textbf{Предметом исследования} являются типизированные промежуточные представления и алгоритмы детерминированного контроля генерации системно-динамических моделей на основе текстовых описаний.

\section*{Научная новизна}

Научная новизна работы заключается в следующем:

\begin{itemize}
    \item Предложен конвейер генерации, в котором LLM формирует краткое описание задачи и план типизированных компонентов, а исполняемая модель строится детерминированно из этого плана.
    \item Разработан метод проверки единиц, топологии потоков и производных в едином каноническом представлении, сопровождаемый ограниченным исправлением и контрольной симуляцией.
    \item Для явно заданных ОДУ реализовано сохранение исходной формы выражений: только верхнеуровневые слагаемые выделяются отдельно, а внешние коэффициенты и скобки не раскрываются бездумно.
    \item Предложена единая схема оценки, включающая исполняемость, структурные и поведенческие свойства, ошибки исправления и вычислительные затраты.
    \item Реализован ограниченный Vensim importer, охватывающий поэлементные массивы, поддиапазоны, исключения и bounded macros без выполнения произвольного кода.
\end{itemize}

Вклад предложенного подхода относительно прямой LLM-генерации оценивается отдельным controlled ablation на заранее зафиксированном corpus. Настоящая работа различает реализованный метод, его инженерную проверку и научную оценку по протоколу, описанному в главе 4. Результаты для конкретной модели и конфигурации inference не обобщаются автоматически на другие checkpoint или произвольные предметные области.

\section*{Практическая значимость}

Созданная платформа позволяет исследовать генерацию моделей в воспроизводимой форме: каждый запрос может сопровождаться промежуточными представлениями, параметрами запуска, результатами проверки и симуляции. Она может применяться для быстрого прототипирования, обучения системной динамике и сравнения методов ограниченной генерации. Поддержка локальных и удалённых OpenAI-compatible провайдеров даёт возможность выбирать между приватностью локального вычисления и качеством более крупных моделей. При этом высокие требования к качеству аннотации и проверки не позволяют рассматривать систему как замену эксперта.

\section*{Структура работы}

В первой главе рассматриваются основы системной динамики, существующие подходы к LLM-генерации и требования к исследуемому методу. Вторая глава описывает типизированное представление, детерминированные построители, проверку и ограниченное исправление. Третья глава посвящена реализации модулей, пользовательских сценариев и вычислительной инфраструктуры. Четвёртая глава содержит результаты экспериментов, анализ ошибок, ограничения и протокол дальнейшей независимой оценки. В приложениях приведены конфигурация, правила оценки, иллюстративный пример и команды воспроизведения.
